\documentclass[../cours_1re.tex]{subfiles}
\begin{document}
\setcounter{chapter}{16}
\chapter{Longueurs, angles et lieux géométriques}
\label{ch:17}
\begin{objectif}
Démontrer les relations métriques et les lieux géométriques ; choisir la représentation la plus courte.
Prérequis : produit scalaire, milieu et coordonnées de seconde.
\end{objectif}
\section{Le théorème d'Al-Kashi}
\begin{theoreme}
Dans un triangle $ABC$, en notant $a=BC$, $b=CA$, $c=AB$,
\[
a^2=b^2+c^2-2bc\cos\widehat{BAC}.
\]
\end{theoreme}
\begin{demonstration}
$\overrightarrow{BC}=\overrightarrow{AC}-\overrightarrow{AB}$.
La formule du carré d'une différence vectorielle donne
$BC^2=AC^2+AB^2-2\overrightarrow{AC}\cdot\overrightarrow{AB}$ ;
la définition par le cosinus termine le calcul.
\end{demonstration}
Pythagore est le cas $\widehat{BAC}=\pi/2$. Réciproquement,
$a^2=b^2+c^2$ force le cosinus à être nul, donc l'angle à être droit.
Si l'on connaît trois côtés d'un triangle non aplati,
$\cos\widehat{BAC}=(b^2+c^2-a^2)/(2bc)$.

\begin{exemple}
Avec $AB=5$, $AC=7$ et $\widehat{BAC}=\pi/3$,
$BC^2=25+49-35=39$, donc $BC=\sqrt{39}$.
Les longueurs doivent être dans une même unité.
\end{exemple}

\section{Aires et loi des sinus}
La hauteur relative au côté $AB$ mesure $AC\sin\widehat{BAC}$,
y compris pour un angle obtus (on projette sur la droite).
Ainsi $\aire=\frac12bc\sin A$ pour un triangle non aplati.
Les trois expressions de l'aire donnent
\[
\frac{\sin A}{a}=\frac{\sin B}{b}=\frac{\sin C}{c}
=\frac{2\aire}{abc}.
\]
\approfondissement
Si $R$ est le rayon du cercle circonscrit, une corde de longueur $a$
interceptée par l'angle inscrit $A$ vérifie $a=2R\sin A$.
Voici une preuve qui ne suppose pas le théorème de l'angle inscrit.
Plaçons $A=(0,0)$, $B=(c,0)$ et $C=(b\cos A,b\sin A)$,
en choisissant $C$ au-dessus de l'axe. Le centre $\Omega$ du cercle a
pour abscisse $c/2$, car $\Omega A=\Omega B$. Écrivons $\Omega=(c/2,t)$.
L'égalité $\Omega A^2=\Omega C^2$ donne
$cb\cos A+2tb\sin A=b^2$, d'où
$t=(b-c\cos A)/(2\sin A)$. Alors
\[
4R^2\sin^2 A
=c^2\sin^2 A+(b-c\cos A)^2
=b^2+c^2-2bc\cos A=a^2.
\]
Comme $\sin A>0$, prendre les racines donne $2R\sin A=a$.
Permuter les sommets fournit $a/\sin A=b/\sin B=c/\sin C=2R$.

\section{La méthode du milieu}
Soit $I$ le milieu de $[AB]$. Pour tout point $M$ :
\[
\overrightarrow{MA}\cdot\overrightarrow{MB}
=MI^2-\frac{AB^2}{4},\qquad
MA^2+MB^2=2MI^2+\frac{AB^2}{2}.
\]
\begin{demonstration}
Posons $\vec u=\overrightarrow{MI}$ et $\vec v=\overrightarrow{IA}$.
Comme $\overrightarrow{IB}=-\vec v$,
$\overrightarrow{MA}=\vec u+\vec v$ et
$\overrightarrow{MB}=\vec u-\vec v$.
Le produit se développe en $\|\vec u\|^2-\|\vec v\|^2$.
La somme des carrés des normes annule les termes croisés.
Enfin $IA=AB/2$.
\end{demonstration}
\begin{theoreme}[cercle de diamètre]
Si $A\ne B$, l'ensemble des points $M$ tels que
$\overrightarrow{MA}\cdot\overrightarrow{MB}=0$ est le cercle
de centre $I$ et de rayon $AB/2$, c'est-à-dire de diamètre $[AB]$.
\end{theoreme}
\begin{demonstration}
La formule précédente transforme l'égalité en $MI^2=AB^2/4$,
équivalente à $MI=AB/2$. Les étapes sont réversibles.
Les points $A,B$ appartiennent au lieu ; pour eux on ne parle pas d'angle
$AMB$, puisque l'un des vecteurs est nul.
\end{demonstration}
Si $A=B$, le lieu est le seul point $A$.
Plus généralement, l'équation
$\overrightarrow{MA}\cdot\overrightarrow{MB}=k$
donne un cercle, un point ou l'ensemble vide selon le signe de $k+AB^2/4$.

\section{Pourquoi les médianes sont concourantes}
\approfondissement
Dans un repère, définissons $G=(A+B+C)/3$ : cette notation représente
la moyenne des coordonnées des trois points. Si $I=(B+C)/2$,
$G=A+\frac23(I-A)$. Ainsi $G$ est sur la médiane issue de $A$,
aux deux tiers depuis le sommet. Le même calcul aux autres sommets prouve
la concurrence des trois médianes. La définition est indépendante du repère :
$\overrightarrow{GA}+\overrightarrow{GB}+\overrightarrow{GC}=\vec0$.
Pour tout $M$, le développement des carrés donne
\[
MA^2+MB^2+MC^2=3MG^2+GA^2+GB^2+GC^2.
\]
Les termes croisés s'annulent grâce à la somme vectorielle nulle.
$G$ minimise donc la somme des trois distances au carré.

\section{Pourquoi les hauteurs sont concourantes}
\approfondissement
Dans un triangle non aplati, les hauteurs issues de $A$ et de $B$ se coupent
en un point $H$ : leurs directions ne sont pas parallèles puisque
$(BC)$ et $(AC)$ ne le sont pas. Or, pour tout point $H$,
\[
\overrightarrow{HA}\cdot\overrightarrow{BC}
+\overrightarrow{HB}\cdot\overrightarrow{CA}
+\overrightarrow{HC}\cdot\overrightarrow{AB}=0.
\]
Pour le vérifier, poser $\vec a=\overrightarrow{HA}$,
$\vec b=\overrightarrow{HB}$, $\vec c=\overrightarrow{HC}$ :
l'expression devient
$\vec a\cdot(\vec c-\vec b)+\vec b\cdot(\vec a-\vec c)
+\vec c\cdot(\vec b-\vec a)$, dont les termes s'annulent deux à deux.
Les deux premières expressions sont nulles par définition de $H$ ;
la troisième est donc nulle. $H$ est aussi sur la troisième hauteur.

\begin{finchapitre}
\section{Utilité et prolongements}
Une équation vectorielle révèle souvent le bon centre d'un problème.
La méthode du milieu et celle du centre de gravité annoncent la décomposition
« dispersion autour du centre + distance au centre ».

Cette même identité réapparaît pour la variance au chapitre 22 et pour
les moindres carrés au chapitre 27. Les barycentres seront généralisés
aux masses, aux intégrales et aux probabilités.
\begin{retenir}
Al-Kashi généralise Pythagore. Pour un lieu géométrique, établir
une équivalence et traiter les cas vide, point et cercle. Les concours de droites se démontrent.
\end{retenir}
\end{finchapitre}
\end{document}
